\documentclass[10pt,a4paper]{amsart}
\usepackage{amsmath,amssymb,amsthm,graphicx}
% For commutative diagrams we use Xy-pic
\usepackage[all]{xy}

% We use 2cell for 2-commutative diagrams.
\xyoption{2cell}
% Theorem environments.
%
\theoremstyle{plain}
\newtheorem{theorem}[subsection]{Theorem}
\newtheorem{proposition}[subsection]{Proposition}
\newtheorem{lemma}[subsection]{Lemma}

\theoremstyle{definition}
\newtheorem{definition}[subsection]{Definition}
\newtheorem{example}[subsection]{Example}
\newtheorem{exercise}[subsection]{Exercise}
\newtheorem{situation}[subsection]{Situation}

\theoremstyle{remark}
\newtheorem{remark}[subsection]{Remark}
\newtheorem{remarks}[subsection]{Remarks}

\numberwithin{equation}{subsection}

% Macros
%
\def\lim{\mathop{\mathrm{lim}}\nolimits}
\def\colim{\mathop{\mathrm{colim}}\nolimits}
\def\Spec{\mathop{\mathrm{Spec}}}
\def\Hom{\mathop{\mathrm{Hom}}\nolimits}
\def\Ext{\mathop{\mathrm{Ext}}\nolimits}
\def\SheafHom{\mathop{\mathcal{H}\!\mathit{om}}\nolimits}
\def\SheafExt{\mathop{\mathcal{E}\!\mathit{xt}}\nolimits}
\def\Sch{\mathit{Sch}}
\def\Mor{\mathop{\mathrm{Mor}}\nolimits}
\def\Ob{\mathop{\mathrm{Ob}}\nolimits}
\def\Sh{\mathop{\mathit{Sh}}\nolimits}
\def\NL{\mathop{N\!L}\nolimits}
\def\CH{\mathop{\mathrm{CH}}\nolimits}
\def\proetale{{pro\text{-}\acute{e}tale}}
\def\etale{{\acute{e}tale}}
\def\QCoh{\mathit{QCoh}}
\def\Ker{\mathop{\mathrm{Ker}}}
\def\Im{\mathop{\mathrm{Im}}}
\def\Coker{\mathop{\mathrm{Coker}}}
\def\Coim{\mathop{\mathrm{Coim}}}

% Boxtimes
%
\DeclareMathSymbol{\boxtimes}{\mathbin}{AMSa}{"02}

%
% Macros for moduli stacks/spaces
%
\def\QCohstack{\mathcal{QC}\!\mathit{oh}}
\def\Cohstack{\mathcal{C}\!\mathit{oh}}
\def\Spacesstack{\mathcal{S}\!\mathit{paces}}
\def\Quotfunctor{\mathrm{Quot}}
\def\Hilbfunctor{\mathrm{Hilb}}
\def\Curvesstack{\mathcal{C}\!\mathit{urves}}
\def\Polarizedstack{\mathcal{P}\!\mathit{olarized}}
\def\Complexesstack{\mathcal{C}\!\mathit{omplexes}}
% \Pic is the operator that assigns to X its picard group, usage \Pic(X)
% \Picardstack_{X/B} denotes the Picard stack of X over B
% \Picardfunctor_{X/B} denotes the Picard functor of X over B
\def\Pic{\mathop{\mathrm{Pic}}\nolimits}
\def\Picardstack{\mathcal{P}\!\mathit{ic}}
\def\Picardfunctor{\mathrm{Pic}}
\def\Deformationcategory{\mathcal{D}\!\mathit{ef}}
\usepackage{cleveref}
\crefname{originallemma}{original lemma}{original lemmas}
\Crefname{originallemma}{Original lemma}{Original lemmas}
\makeatletter
\newlabel{lemma-from-chow-to-K}{{0FDS}{source}{}{sourcepointer-1}{}}
\newlabel{lemma-from-chow-to-K@cref}{{[originallemma][0][]0FDS}{[1][0][]source}{}{}{}}
\makeatother
\title{K\_0 of the coherent sheaves set-theoretically supported on a closed subscheme is K\_0 of the closed subscheme}
\author{The Stacks Project Authors}
\date{stacks-project:chow.tex:local-snapshot-20261012; GNU FDL 1.2 or later}
\begin{document}
\maketitle
\noindent\textit{Editorial note.} All delivered mathematical text is the author's own, copied from the pinned Stacks Project chapter file chow.tex; no statement or proof step was rewritten, reordered or removed.\par\medskip
\noindent\textit{Editorial note.} Preamble provenance: the pinned chapter file opens with \textbackslash{}input\{preamble\} and carries no preamble text of its own, so the project's single Stacks-wide preamble staged at sources/stacks-chow/preamble.tex is excerpted verbatim as two spans, lines 14-18 (Xy-pic setup) and lines 154-217 (the native theorem declarations and mathematical macros). Nothing was added, removed or redefined; the wrapper class and the amsmath/amssymb/amsthm/graphicx base packages are supplied by the assembler.\par\medskip
\noindent\textit{Editorial note.} Same-chapter pointers are located at their \textbackslash{}label lines inside the pinned chapter file; cross-chapter pointers are located in the project's own tag list staged in the snapshot (sources/stacks-chow/tags.txt), at the verbatim lines carrying their canonical tags.\par\medskip
\noindent\textit{Editorial note.} The author's closing sentence at source line 4125, 'Some details omitted.', is retained verbatim at the end of the counted proof. The filtration, the construction of the inverse map and the double-filtration comparison that establishes additivity are all carried out inside the span; only repeated re-derivations of the same local-finiteness argument are left to the reader.\par\medskip
\noindent\textit{Allowed standard prerequisites (editorial).}\begin{itemize}
\item The locally finite I-adic filtration and the replacement of a class by its graded pieces. Same chapter chow.tex, label lemma-from-chow-to-K, invoked at proof lines 4071 and 4083. Supplies the shape of the argument: the descending filtration by powers of the ideal sheaf is locally finite, so the direct sums of the F\textasciicircum{}p and of the graded pieces are coherent and the class of a module may be replaced by the class of its graded pieces. The proof here re-runs that argument inside the subcategory supported on Z and then goes on to construct the inverse map and prove its additivity, which the cited lemma does not do.
\end{itemize}\medskip
\ifdefined\hypertarget\hypertarget{sourcepointer-1}{}\fi
\noindent\textit{Source reference.} Original lemma 0FDS: sources/stacks-chow/pinned-original-source.tex, source lines 3934--3934. Navigation pointer for \textbackslash{}ref\{lemma-from-chow-to-K\} used inside the counted span (line 4071); the argument that consumes it is retained in the span.\par\medskip
\begin{lemma}
\label{lemma-K-coherent-supported-on-closed}
Let $X$ be a locally Noetherian scheme. Let $Z \subset X$ be a closed
subscheme. Denote $\textit{Coh}_Z(X) \subset \textit{Coh}(X)$
the Serre subcategory of coherent $\mathcal{O}_X$-modules whose
set theoretic support is contained in $Z$. Then the exact inclusion
functor $\textit{Coh}(Z) \to \textit{Coh}_Z(X)$ induces
an isomorphism
$$
K'_0(Z) = K_0(\textit{Coh}(Z)) \longrightarrow K_0(\textit{Coh}_Z(X))
$$
\end{lemma}

\begin{proof}
Let $\mathcal{F}$ be an object of $\textit{Coh}_Z(X)$.
Let $\mathcal{I} \subset \mathcal{O}_X$ be the quasi-coherent
ideal sheaf of $Z$. Consider the descending filtration
$$
\ldots \subset
\mathcal{F}^p = \mathcal{I}^p \mathcal{F} \subset
\mathcal{F}^{p - 1} \subset \ldots \subset \mathcal{F}^0 = \mathcal{F}
$$
Exactly as in the proof of Lemma \ref{lemma-from-chow-to-K} this filtration
is locally finite and hence
$\bigoplus_{p \geq 0} \mathcal{F}^p$,
$\bigoplus_{p \geq 1} \mathcal{F}^p$, and
$\bigoplus_{p \geq 0} \mathcal{F}^p/\mathcal{F}^{p + 1}$
are coherent $\mathcal{O}_X$-modules supported on $Z$.
Hence we get
$$
[\mathcal{F}] =
[\bigoplus\nolimits_{p \geq 0} \mathcal{F}^p/\mathcal{F}^{p + 1}]
$$
in $K_0(\textit{Coh}_Z(X))$ exactly as in the proof of
Lemma \ref{lemma-from-chow-to-K}. Since the coherent module
$\bigoplus_{p \geq 0} \mathcal{F}^p/\mathcal{F}^{p + 1}$
is annihilated by $\mathcal{I}$ we conclude that
$[\mathcal{F}]$ is in the image. Actually, we claim that the map
$$
\mathcal{F} \longmapsto 
c(\mathcal{F}) =
[\bigoplus\nolimits_{p \geq 0} \mathcal{F}^p/\mathcal{F}^{p + 1}]
$$
factors through $K_0(\textit{Coh}_Z(X))$ and is an inverse to
the map in the statement of the lemma. To see this all we have
to show is that if
$$
0 \to \mathcal{F} \to \mathcal{G} \to \mathcal{H} \to 0
$$
is a short exact sequence in $\textit{Coh}_Z(X)$, then we
get $c(\mathcal{G}) = c(\mathcal{F}) + c(\mathcal{H})$.
Observe that for all $q \geq 0$ we have a short exact sequence
$$
0 \to
(\mathcal{F} \cap \mathcal{I}^q\mathcal{G})/
(\mathcal{F} \cap \mathcal{I}^{q + 1}\mathcal{G}) \to
\mathcal{G}^q/\mathcal{G}^{q + 1} \to
\mathcal{H}^q/\mathcal{H}^{q + 1} \to 0
$$
For $p, q \geq 0$ consider the coherent submodule
$$
\mathcal{F}^{p, q} = \mathcal{I}^p\mathcal{F} \cap \mathcal{I}^q\mathcal{G}
$$
Arguing exactly as above and using that the filtrations
$\mathcal{F}^p = \mathcal{I}^p\mathcal{F}$ and
$\mathcal{F} \cap \mathcal{I}^q\mathcal{G}$ are locally finite,
we find that
$$
[\bigoplus\nolimits_{p \geq 0} \mathcal{F}^p/\mathcal{F}^{p + 1}] =
[\bigoplus\nolimits_{p, q \geq 0}
\mathcal{F}^{p, q}/(\mathcal{F}^{p + 1, q} + \mathcal{F}^{p, q + 1})] =
[\bigoplus\nolimits_{q \geq 0}
(\mathcal{F} \cap \mathcal{I}^q\mathcal{G})/
(\mathcal{F} \cap \mathcal{I}^{q + 1}\mathcal{G})]
$$
in $K_0(\textit{Coh}(Z))$. Combined with the exact sequences above we obtain
the desired result. Some details omitted.
\end{proof}

\end{document}
